% Introduction — ReCAHS paper

\section{Introduction}
\label{sec:intro}

Transformer forecasters such as PatchTST \citep{nie2023patchtst} and
iTransformer \citep{liu2024itransformer} dominate long-term time-series
benchmarks, yet much of their multi-head attention is redundant: a
substantial fraction of heads can be removed with little loss
\citep{michel2019sixteen,voita2019analyzing}. The standard recipe scores
each head with one global importance value and prunes the lowest-scoring
ones. Time series, however, are not homogeneous: windows alternate between
trend-dominant, seasonal-dominant and noise-dominant \emph{regimes}, and a
head that is useful in one regime can be harmful in another. This suggests
a regime-conditional policy --- detect the regime of each input window and
activate the head subset selected for that regime --- an instance of
input-conditional computation \citep{han2021dynamic} applied to head
sparsity.

Prior comparisons of dynamic against static pruning share a structural
weakness: the dynamic method and its static baseline differ
simultaneously in \emph{how} masks are selected and in \emph{whether}
they adapt to the input, so the observed gains cannot be attributed to
adaptivity. We conduct, to our knowledge, the first controlled study that
separates these factors, on a grid of 4~datasets $\times$ 4~horizons
$\times$ 5~seeds $\times$ 6~sparsity levels (${\sim}9{,}300$ masked
evaluations), with per-window losses stored for every configuration so
that all statistics are recomputable offline.

Our contributions:
\begin{enumerate}
  \item \textbf{Fragility of the standard method.} Global importance
  ranking degrades sharply at realistic sparsity and is beaten by
  \emph{random} pruning on minute-resolution data; greedy backward
  elimination (joint selection) stays within ${\sim}1\%$ of the unpruned
  model at $25\%$ sparsity and degrades gracefully to $75\%$
  (Sec.~\ref{sec:results-sweep}).
  \item \textbf{Controlled attribution.} Regime-routed selection beats
  the standard static baseline in up to $15/16$ cells (sign test
  $p{=}5\times10^{-4}$), but a pooled greedy mask --- same criterion, no
  regime conditioning --- reverses the comparison: the criterion, not the
  routing, explains the gain (Sec.~\ref{sec:results-attribution}).
  \item \textbf{The routing headroom is real but aleatoric.} A per-window
  oracle over the regime masks beats even the \emph{unpruned} model in
  $15/16$ cells (up to $-12\%$), yet twelve deployable routers ---
  STL-based, learned, bandit, gated, representation-based --- fail to
  capture it. Near-oracle in-sample fits, flat $10\times$ learning curves
  and failed zero-shift diagnostics locate the failure in the signal, not
  the router (Sec.~\ref{sec:results-routing}).
  \item \textbf{Combination instead of selection.} Averaging the
  specialist masks' predictions beats the best static mask on all four
  datasets and surpasses the unpruned model on two; a soft-routing
  $\alpha$-sweep shows performance \emph{increases} as trust in the
  regime label decreases. An ablation attributes the gain to mask
  diversity, for which regime partitioning is the strongest but not a
  necessary source (Sec.~\ref{sec:results-ensemble}).
  \item \textbf{Practical results.} With the mask imposed from scratch
  during training, $25\%$ head pruning is free ($-5.4\%$ parameters,
  $-12\%$ FLOPs); a $12$-feature classifier replaces per-window STL
  detection at $45\times$ lower latency with no end-to-end accuracy loss;
  and the findings transfer with a precise boundary to iTransformer's
  cross-variate attention (Sec.~\ref{sec:results-itr}).
\end{enumerate}
